% Part:sets-functions-relations
% Chapter: size-of-sets
% Section: pairing

\documentclass[../../../include/open-logic-section]{subfiles}

\begin{document}

\olfileid[fa-IR]{sfr}{siz}{pai}
\olsection{توابع جفت‌سازی و کدها}

\begin{explain}
روش زیگزاگ کانتور شماراییِ $\Nat^n$ را به‌صورت دیداری آشکار می‌کند. اما
بیایید بر آرایه‌ای متمرکز شویم که $\Nat^2$ را نمایش می‌دهد. با دنبال‌کردن
خط زیگزاگ در آرایه و شمردن جایگاه‌ها، می‌توانیم بررسی کنیم که
$\tuple{1,2}$ با عدد~$7$ متناظر است. بااین‌حال، خوب بود اگر می‌توانستیم
این مقدار را مستقیم‌تر محاسبه کنیم. یعنی خوب بود \emph{وارونِ} برشماری
زیگزاگ، یعنی $g\colon \Nat^2 \to \Nat$، را در دسترس داشتیم، به‌گونه‌ای که
\[
g(\tuple{0,0}) = 0, \;
g(\tuple{0,1}) = 1, \;
g(\tuple{1,0}) = 2, \; \dots,
g(\tuple{1,2}) = 7, \; \dots
\]
این تابع به ما امکان می‌داد دقیقاً محاسبه کنیم که $\tuple{n, m}$ در کجای
برشماری ظاهر خواهد شد.

در واقع، می‌توانیم با دو مشاهده $g$ را مستقیماً تعریف کنیم. نخست: اگر
سطر $n$اُم و ستون $m$اُم حاوی مقدار~$v$ باشد، آنگاه سطر
$(n+1)$اُم و ستون $(m-1)$اُم حاوی مقدار $v + 1$ است. دوم: سطر نخست
برشماری ما از اعداد مثلثی تشکیل می‌شود که با $0$، $1$، $3$، $6$ و
الی آخر آغاز می‌شوند. عدد مثلثیِ $k$اُم مجموع اعداد طبیعیِ $< k$ است
که می‌توان آن را به‌صورت $k(k+1)/2$ محاسبه کرد. با کنار هم گذاشتن این
دو مشاهده، تابع زیر را در نظر بگیرید:
\[
  g(n,m) = \frac{(n+m+1)(n+m)}{2} + n
\]
اغلب صرفاً $g(n, m)$ را به‌جای $g(\tuple{n, m})$ می‌نویسیم، زیرا خواندنش
آسان‌تر است. این فرمول می‌گوید نخست عدد مثلثیِ $(n+m)^\text{اُم}$ را
به دست آورید و سپس $n$ را به آن بیفزایید. این کار آرایه را دقیقاً به
شیوهٔ مطلوب ما پر می‌کند. پس به‌ویژه، زوج $\tuple{1, 2}$ به
$\frac{4 \times 3}{2} +
1 = 7$ فرستاده می‌شود.

این تابع $g$ \emph{وارونِ} برشماریِ مجموعه‌ای از زوج‌هاست. چنین توابعی
\emph{توابع جفت‌سازی} نامیده می‌شوند.
\end{explain}

\begin{defn}[تابع جفت‌سازی]
  تابع $f\colon A \times B \to \Nat$ یک \emph{تابع جفت‌سازیِ} حسابی است
  اگر $f$ یک‌به‌یک باشد. همچنین می‌گوییم $f$ \emph{کدگذاریِ}
  $A \times B$ را انجام می‌دهد و $f(x,y)$ \emph{کدِ} $\tuple{x,y}$ است.
\end{defn}

\begin{explain}
می‌توانیم، برای نمونه، زوج‌های اعداد طبیعی را با توابع جفت‌سازی
کدگذاری کنیم؛ یا به بیان دیگر، هر \emph{زوج} از عناصر را با
\emph{یک} عدد بازنمایی کنیم. با استفاده از وارون تابع جفت‌سازی می‌توانیم
عدد را \emph{کدگشایی کنیم}؛ یعنی دریابیم کدام زوج را بازنمایی می‌کند.
\end{explain}

\begin{prob}
برشماری‌ای از مجموعهٔ همهٔ اعداد گویای نامنفی ارائه کنید.
\end{prob}

\begin{prob}
نشان دهید $\Rat$ !!{enumerable} است. به یاد آورید که هر عدد گویا را
می‌توان به‌صورت کسر $z/m$ با $z \in \Int$ و $m \in \Nat^+$ نوشت.
\end{prob}

\begin{prob}
برشماری‌ای از $\Bin^*$ تعریف کنید.
\end{prob}

\begin{prob}
از درس مقدماتی منطق خود به یاد آورید که هر جدول ارزشِ ممکن یک تابع ارزش
را بیان می‌کند. به بیان دیگر، توابع ارزش همهٔ توابع از
$\Bin^k \to \Bin$ برای یک~$k$ هستند. ثابت کنید مجموعهٔ همهٔ توابع ارزش
شمارا است.
\end{prob}

\begin{prob}
نشان دهید مجموعهٔ همهٔ زیرمجموعه‌های متناهیِ یک مجموعهٔ نامتناهی و
!!{enumerable} دلخواه، !!{enumerable} است.
\end{prob}

\begin{prob}
زیرمجموعه‌ای از $\Nat$ را \emph{هم‌متناهی} می‌نامیم اگر و تنها اگر
متممِ یک زیرمجموعهٔ متناهی از~$\Nat$ باشد؛ یعنی $A \subseteq \Nat$
هم‌متناهی است اگر و تنها اگر $\Nat\setminus A$ متناهی باشد. فرض کنید
$I$ مجموعه‌ای باشد که !!{element}s آن دقیقاً زیرمجموعه‌های متناهی و
هم‌متناهیِ $\Nat$ هستند. نشان دهید $I$ !!{enumerable} است.
\end{prob}

\begin{prob}
نشان دهید اجتماعِ خانواده‌ای !!{enumerable} از مجموعه‌های !!{enumerable}،
!!{enumerable} است. یعنی هرگاه $A_1$، $A_2$، \dots{} مجموعه باشند و هر
$A_i$ !!{enumerable} باشد، آنگاه اجتماع $\bigcup_{i=1}^\infty
A_i$ همهٔ آن‌ها نیز !!{enumerable} است. [توجه: این مسئله دشوار است!]
\end{prob}

\begin{prob}
فرض کنید $f \colon A \times B \to \Nat$ یک تابع جفت‌سازی دلخواه باشد.
نشان دهید وارونِ $f$ یک برشماری از $A \times B$ است.
\end{prob}

\begin{prob}
تابعی مشخص کنید که $\Nat^3$ را کدگذاری کند.
\end{prob}

\end{document}
